% Generated draft: requires mathematical and source review.
\documentclass[11pt]{article}
\usepackage[margin=25mm]{geometry}
\usepackage{fontspec}
\setmainfont{Latin Modern Roman}
\usepackage{amsmath,amssymb,mathtools,unicode-math}
\setmathfont{Latin Modern Math}
\usepackage{enumitem,needspace,xurl,hyperref}
\setlength{\parindent}{0pt}
\setlength{\parskip}{5pt}
\setlength{\emergencystretch}{3em}
\newcommand{\N}{\mathbb{N}}
\newcommand{\Z}{\mathbb{Z}}
\newcommand{\Q}{\mathbb{Q}}
\newcommand{\R}{\mathbb{R}}
\newcommand{\C}{\mathbb{C}}
\newcommand{\F}{\mathbb{F}}
\newcommand{\A}{\mathbb{A}}
\newcommand{\PP}{\mathbb{P}}
\newcommand{\E}{\mathbb{E}}
\newcommand{\eps}{\varepsilon}
\newcommand{\dd}{\,\mathrm{d}}
\begin{document}
\section*{UnsolvedMath ID: 30006194}
\textbf{Source problem code:} catalogue entry
\subsection{Definitions and mathematical statement}
For a connected smooth complex projective variety \(X\), let \(X^{\mathrm{an}}\) denote its underlying complex analytic space, choose a basepoint \(x\in X^{\mathrm{an}}\), and write \(\pi_1(X,x)=\pi_1(X^{\mathrm{an}},x)\) for its topological fundamental group. Its algebraic fundamental group is the profinite group classifying pointed finite étale covers; over \(\C\), it is naturally identified with
\[
\pi_1^{\mathrm{alg}}(X,x)\cong\widehat{\pi_1(X,x)}
=\varprojlim_{\substack{N\trianglelefteq\pi_1(X,x)\\[\pi_1(X,x):N]<\infty}}\pi_1(X,x)/N.
\]
Here isomorphisms of algebraic fundamental groups mean continuous isomorphisms of profinite groups. The isomorphism types in the questions do not depend on the chosen basepoints.

On the connected-variety interpretation of the supplied source, the two reconstruction questions ask whether the following implications hold for every pair of connected smooth complex projective varieties \(X,Y\) and choices of basepoints \(x,y\):
\begin{enumerate}
\item Does
\[
\pi_1^{\mathrm{alg}}(X,x)\cong\pi_1^{\mathrm{alg}}(Y,y)
\quad\Longrightarrow\quad
X^{\mathrm{an}}\cong_{\mathrm{homeo}}Y^{\mathrm{an}}
\]
hold, where the conclusion means that the underlying topological spaces are homeomorphic?
\item Does
\[
\pi_1^{\mathrm{alg}}(X,x)\cong\pi_1^{\mathrm{alg}}(Y,y)
\quad\Longrightarrow\quad
\pi_1(X,x)\cong\pi_1(Y,y)
\]
hold, where the conclusion is an isomorphism of abstract groups?
\end{enumerate}

\textit{Unresolved scope qualifications.} The source does not explicitly state connectedness; the formulation above uses the connected interpretation conventionally needed for a single fundamental group to describe the variety. If disconnected varieties are intended, the treatment of their components must be specified. The source also gives no fixed dimension, asphericity hypothesis, or restricted comparison family. These implications transcribe the unrestricted reconstruction questions on the connected interpretation; whether the original discussion intended additional restrictions requires review.

\textit{Supplied literature qualification.} The registry labels the problem partially solved. Its dated assessment reports profinite rigidity and homeomorphism recovery for important families, including products of hyperbolic Riemann surfaces, and describes both questions for arbitrary smooth complex projective varieties as remaining open. This is supplied curation, not an independently verified mathematical or literature conclusion. The record does not supply precise statements or comparison classes for those reported results.
\subsection{Short English statement}
The source asks whether the algebraic fundamental group determines a smooth complex projective variety up to homeomorphism, and whether it determines the variety's topological fundamental group up to group isomorphism. The displayed formulation interprets determination by comparison with every other connected smooth complex projective variety. Connectedness and any intended restrictions on the comparison class require review. The supplied partially solved status and literature assessment are dated curation, not independently verified conclusions.
\subsection{Sources}
\begin{itemize}
\item Local UnsolvedMath registry, permanent ID \(30006194\), title \emph{Profinite Reconstruction of Complex Projective Varieties}, exact source problem code \texttt{catalogue entry}; supplied complete source record in the Oberwolfach Reports Open Problems collection, proposed year 2025.
\item Supplied report citation: \emph{Geometric Structures in Group Theory} (2025), Oberwolfach Reports, DOI 10.4171/owr/2025/6, pp.~275--338. \url{https://doi.org/10.4171/owr/2025/6}. The supplied report citation gives a report-wide page range rather than the exact problem page.
\item Attribution in the supplied original statement: Claudio Llosa Isenrich, \emph{Profinite rigidity of Kähler groups}, joint work with Sam Hughes, Pierre Py, Matthew Stover, and Stefano Vidussi. The registry's separate proposer field is empty; this attribution is to the supplied discussion and joint work.
\item Supplied literature source: Hughes et al., profinite rigidity of Kähler groups. \url{https://arxiv.org/abs/2501.13761}.
\item Supplied literature source: OWR 2025/6 geometric group theory report. \url{https://oa.tib.eu/renate/server/api/core/bitstreams/6f61d540-8e50-41f3-97a8-6e8214672460/content}.
\item Literature assessment date: \texttt{2026-08-22}. This is supplied dated curation, not independently verified. The supplied assessment reports partial results for important families such as products of hyperbolic Riemann surfaces and labels the unrestricted questions open; its time-sensitive status and cited claims require independent verification.
\item The registry labels the statement \texttt{corrected\_verified} and reports that it was reconstructed and checked against the official Oberwolfach report. These are supplied verification claims; no independent source inspection is asserted here.
\end{itemize}
\end{document}
